% !TeX root = ../../../金融经济学_王江_讲义及习题.tex
% 同一纯风险前沿：\sigma^2=1+((\bar r-3.2)/.7)^2，故 c=1。
% 左图 r_F=3.2，K=.49。右图 r_F=2.65，H=.55，K=.49+.55^2。
\begin{tikzpicture}[
 >=Stealth,font=\small,x=1.48cm,y=.66cm,
 risky/.style={red!65!black,thick},
 upper/.style={blue!70!black,very thick},
 lower/.style={blue!70!black,thick,dashed},
 guide/.style={gray,densely dotted}]
 \pgfmathsetmacro{\SlopeT}{sqrt(.49+.55^2)}
 \pgfmathsetmacro{\SigmaT}{\SlopeT/.55}
 \pgfmathsetmacro{\MeanT}{2.65+(.49+.55^2)/.55}

 \foreach \Shift/\RiskFree/\Slope/\Panel in
   {0/3.2/.7/{(a) $H=0$：无切点},6.4/2.65/\SlopeT/{(b) $H>0$：在上支相切}} {
   \begin{scope}[xshift=\Shift cm]
     \draw[->] (0,-.35)--(3.35,-.35) node[right] {$\sigma_p$};
     \draw[->] (0,-.35)--(0,6.03) node[above] {$\bar r_p$};
     \node[below left,inner sep=2pt] at (0,-.35) {$0$};
     \draw[guide] (0,3.2)--(1,3.2);
     \draw[risky] plot[domain=-2.93:2.93,samples=121,variable=\t]
       ({sqrt(1+\t*\t)},{3.2+.7*\t});
     \draw[upper] (0,\RiskFree)--(3.12,{\RiskFree+3.12*\Slope});
     \draw[lower] (0,\RiskFree)--(3.12,{\RiskFree-3.12*\Slope});
     \fill (0,\RiskFree) circle (1.7pt) node[left] {$F$};
     \fill (1,3.2) circle (1.7pt) node[right=3pt] {纯风险 MVP};
     \node at (1.65,-.92) {\Panel};
   \end{scope}
 }
 \begin{scope}[xshift=6.4cm]
   \fill (0,3.2) circle (.9pt);
   \fill[blue!70!black] (\SigmaT,\MeanT) circle (2pt) node[above left=1pt] {$T$};
 \end{scope}

 \draw[risky] (.05,-1.65)--(.43,-1.65) node[right=2pt,black] {纯风险前沿};
 \draw[upper] (2.35,-1.65)--(2.73,-1.65) node[right=2pt,black] {含无风险证券：上支};
 \draw[lower] (5.6,-1.65)--(5.98,-1.65) node[right=2pt,black] {下支};
\end{tikzpicture}
